\chapter{WP2: Soluzione}
\label{ch:wp2}

% -----------------------------------------------------------------------------
% Nota strutturale
% -----------------------------------------------------------------------------
% Questo file mantiene una struttura da completare in LaTeX, ma aggiornata alla
% versione WP2_Soluzione (2).md. Le modifiche principali rispetto alla struttura
% precedente sono:
% - introduzione del commitment pubblico K_j = H(C_j);
% - ridefinizione del fingerprint come F_j = H(K_j || eta_j);
% - distinzione esplicita tra archivio interno AE e urna pubblica;
% - terminologia "assenza di ricevuta rivelante" al posto di receipt-freeness piena;
% - ridimensionamento del ruolo di TLS rispetto alla traffic analysis;
% - pubblicazione per finestre dei record dell'urna;
% - chiarimento che le firme RSA non sono textbook RSA.


Il protocollo proposto segue un'architettura client-server crittografica e si articola in cinque fasi principali: setup, autenticazione, costruzione della scheda, conservazione nell'urna e scrutinio. L'obiettivo è realizzare una procedura di voto coerente con il modello definito in WP1, mantenendo separati i compiti di identificazione dell'elettore, raccolta delle schede e pubblicazione dei risultati. 

La scelta centrale dell'architettura è la separazione tra Sistema di Autenticazione, indicato come SA, e Autorità Elettorale, indicata come AE. Il SA verifica l'identità degli elettori e controlla che ciascuno riceva una sola abilitazione al voto. L'AE, invece, riceve e conserva le schede cifrate, registra i dati pubblici dell'urna e svolge lo scrutinio finale. In questo modo il SA conosce le identità reali, ma non vede le preferenze espresse; l'AE gestisce le schede e gli pseudonimi crittografici, ma non deve conoscere l'identità reale degli elettori. 

Lo pseudo-anonimato ottenuto dal protocollo dipende quindi dall'Assunzione~1 del WP1, cioè dalla non-collusione tra SA e AE. Se le due entità condividessero i propri registri interni, il legame tra identità reale e scheda potrebbe essere ricostruito. Questa assunzione non viene nascosta, ma diventa uno dei punti di fiducia espliciti del sistema. 

Restano inoltre due limiti strutturali che il protocollo non può eliminare usando solo gli strumenti trattati nel corso. Il primo riguarda la segretezza delle singole schede rispetto all'AE durante lo scrutinio, discussa nella Sezione~\ref{subsec:wp2_limite_sec1}. Il secondo riguarda la verificabilità universale piena del conteggio, discussa nella Sezione~\ref{subsec:wp2_limite_ver2}. In entrambi i casi, WP2 descrive le mitigazioni adottate, mentre WP3 analizzerà il rischio residuo rispetto al threat model.

\clearpage
\subsection*{Requisito TLS e limiti rispetto alla traffic analysis}
\label{subsec:wp2_tls}
% TODO: Dichiarare che tutte le comunicazioni Elettore--SA ed Elettore--AE avvengono
% su TLS con autenticazione del server tramite certificato emesso dalla CA d'Ateneo.
%
% TODO: Motivare TLS come requisito strutturale: la firma applicativa protegge
% l'integrità del pacchetto M, ma non protegge la confidenzialità dei campi trasmessi
% nel pacchetto, tra cui eta, C_p e sigma.
%
% TODO: Specificare che un avversario di rete che osservasse eta potrebbe calcolare
% H(K || eta), con K = H(C), e cercare il fingerprint F_j nell'urna pubblica.
%
% TODO: Correggere esplicitamente il possibile overclaim: TLS protegge i contenuti,
% ma non elimina completamente timing, indirizzi, dimensioni e pattern del traffico.
%
% TODO: Collegare questa mitigazione a due scelte applicative: t_finestra nei
% certificati pseudonimi e pubblicazione per finestre dei record dell'urna.

Tutte le comunicazioni di rete tra l'Elettore e le componenti server (Sistema di Autenticazione e Autorità Elettorale) devono avvenire esclusivamente su canali cifrati tramite protocollo TLS (Transport Layer Security). In questa fase, l'autenticazione dei server è garantita dall'utilizzo di certificati digitali X.509 emessi dalla Certificate Authority (CA) d'Ateneo.

L'impiego di TLS non rappresenta un semplice dettaglio implementativo, ma un requisito strutturale indispensabile. Come verrà dettagliato in seguito, il protocollo applicativo prevede che l'elettore firmi il pacchetto di voto $M$ per garantirne l'integrità. Tuttavia, tale firma non offre alcuna garanzia di confidenzialità. Senza la cifratura del canale di trasporto, i metadati essenziali contenuti nel pacchetto — tra cui il certificato pseudonimo $C_p$, il nonce $\eta$ e la firma $\sigma$ — viaggerebbero in chiaro. 

Questa esposizione vanificherebbe il disaccoppiamento logico tra elettore e scheda. Un avversario di rete passivo, intercettando il traffico in chiaro, potrebbe facilmente estrarre il ciphertext $C$ e il nonce $\eta$. A quel punto, l'attaccante potrebbe calcolare autonomamente il commitment $K = H(C)$ e il conseguente fingerprint $F_j = H(K \parallel \eta)$. Cercando il valore $F_j$ all'interno dell'urna pubblica, l'avversario riuscirebbe a tracciare il voto, associando il record pubblicato all'indirizzo IP del mittente, violando così il requisito di segretezza.

Tuttavia, è fondamentale non sovrastimare le reali garanzie di sicurezza offerte dal canale cifrato. Sebbene il protocollo TLS protegga il contenuto del payload, esso non è progettato per neutralizzare l'analisi del traffico. Un avversario in ascolto può ancora raccogliere metadati di rete cruciali, quali gli indirizzi IP sorgente e destinazione, le dimensioni dei pacchetti e, soprattutto, il \textit{timing} delle connessioni. 

Un'analisi temporale che correli l'istante in cui un elettore si autentica presso il SA con l'istante in cui l'AE riceve una scheda cifrata potrebbe permettere di ricostruire il legame tra identità e voto, eludendo le protezioni crittografiche. Per mitigare questa specifica vulnerabilità a livello applicativo, il protocollo adotta due strategie complementari:
\begin{enumerate}
    \item \textbf{Timestamp a bassa risoluzione:} Il certificato pseudonimo $C_p$ non contiene l'istante esatto di emissione, ma un $t_{finestra}$ che rappresenta un intervallo orario predefinito, riducendo volutamente la precisione del dato temporale a disposizione dell'AE.
    \item \textbf{Pubblicazione per finestre:} L'AE non registra né pubblica i record nell'urna in tempo reale al momento della ricezione, ma li raggruppa e li pubblica a blocchi in finestre temporali sfalsate o randomizzate.
\end{enumerate}
Queste due scelte applicative spezzano la correlazione temporale delle operazioni, ridimensionando fortemente il rischio di de-anonimizzazione basato sull'analisi del traffico di rete.

\clearpage
\section{Parametri e notazione}
\label{sec:parametri_notazione}
% Obiettivo della sezione: fornire una legenda unica, da usare in tutto WP2.
%
% TODO: Inserire una tabella con tre colonne:
% Simbolo | Significato | Contesto operativo.
%
% TODO: Includere almeno i seguenti simboli:
% N, (PU_ele, PR_ele), (PU_SA, PR_SA), (PU_AE, PR_AE), (PU_v, PR_v),
% C_p, i, r, C, eta, id_ele, t_finestra, t_apertura, t_chiusura, H(.),
% K_j = H(C_j), F_j = H(K_j || eta_j), sigma, R.
%
% TODO: Specificare che sigma e R sono firme digitali basate su hash-and-sign,
% non su RSA textbook.


Al fine di garantire chiarezza espositiva ed evitare ambiguità descrittive nelle sezioni successive, la Tabella \ref{tab:parametri_notazione} definisce in un'unica legenda la notazione crittografica e i parametri operativi utilizzati in tutte le fasi del protocollo.

\begin{table}[h!]
\centering
\renewcommand{\arraystretch}{1.3} % Aumenta leggermente lo spazio verticale per migliore leggibilità
\begin{tabular}{|l|p{4.5cm}|p{7.5cm}|}
\hline
\textbf{Simbolo} & \textbf{Significato} & \textbf{Contesto operativo} \\ \hline
$N$ & Numero di opzioni di voto valide. & Utilizzato per il controllo di conformità strutturale della scheda ($1 \le i \le N$) in fase di scrutinio. \\ \hline
$id_{elezione}$ & Identificatore univoco della consultazione in corso. & Previene attacchi di \textit{replay cross-elezione}. \\ \hline
$(PU_{ele}, PR_{ele})$ & Coppia RSA dell'elezione. & $PU_{ele}$ cifra i voti (lato elettore); $PR_{ele}$ decifra le schede (lato AE in fase di scrutinio). \\ \hline
$(PU_{SA}, PR_{SA})$ & Coppia RSA del SA. & $PR_{SA}$ firma i certificati pseudonimi degli elettori; $PU_{SA}$ ne verifica la validità. \\ \hline
$(PU_{AE}, PR_{AE})$ & Coppia RSA dell'AE. & $PR_{AE}$ firma i record dell'urna, le ricevute e il documento finale; $PU_{AE}$ abilita la verifica pubblica. \\ \hline
$(PU_v, PR_v)$ & Coppia RSA effimera. & Generata localmente dall'elettore, è valida solo per la singola sessione di voto e successivamente distrutta. \\ \hline
$C_p$ & Certificato pseudonimo. & Emesso dal SA, attesta l'abilitazione al voto associandola alla chiave pubblica effimera $PU_v$. \\ \hline
$i$ & Preferenza di voto. & Valore in chiaro scelto dall'elettore e codificato come intero scalare. \\ \hline
$r$ & Fattore di randomizzazione. & Generato dall'elettore per lo schema RSA-OAEP, assicura l'indistinguibilità del ciphertext. \\ \hline
$C$ & Scheda cifrata. & Calcolata dall'elettore come $\text{RSA-OAEP-Enc}(PU_{ele}, i; r)$. \\ \hline
$\eta$ & Nonce a 256 bit. & Generato dall'elettore, inserito in chiaro nel pacchetto per garantire l'anti-replay e il fingerprinting. \\ \hline
$t_{finestra}$ & Finestra temporale. & Identificatore a bassa risoluzione temporale inserito in $C_p$ per mitigare l'analisi del traffico. \\ \hline
$t_{apertura}, t_{chiusura}$ & Timestamp ufficiali. & Definiscono i limiti temporali di validità per l'accettazione delle schede. \\ \hline
$H(\cdot)$ & Funzione hash (SHA-256). & Primitiva utilizzata per \textit{commitment}, \textit{fingerprint} e firme digitali. \\ \hline
$K_j = H(C_j)$ & \textit{Commitment} pubblico. & Traccia del ciphertext $C_j$ inserita nell'urna (impegna l'AE senza esporre il ciphertext in chiaro). \\ \hline
$F_j = H(K_j \parallel \eta_j)$ & \textit{Fingerprint} (impronta). & Impronta indicizzata nel Merkle Tree dell'urna per abilitare la verifica individuale. \\ \hline
$\sigma$ & Firma sul pacchetto. & Apposta dall'elettore con $PR_v$ per garantire l'integrità del pacchetto e la paternità effimera. \\ \hline
$R$ & Ricevuta crittografica. & Rilasciata dall'AE e firmata con $PR_{AE}$ per attestare l'avvenuta ricezione e accettazione del pacchetto. \\ \hline
\end{tabular}
\caption{Legenda dei parametri e della notazione crittografica del protocollo.}
\label{tab:parametri_notazione}
\end{table}

\textbf{Nota fondamentale sulle firme digitali:} Per garantire la totale aderenza ai requisiti di sicurezza e prevenire le vulnerabilità intrinseche in presenza di messaggi scelti dall'avversario, si specifica esplicitamente che le firme digitali del protocollo (come $\sigma$ ed $R$) sono implementate unicamente secondo lo schema \textbf{hash-and-sign}. I dati oggetto della firma vengono sempre preventivamente compressi tramite la funzione hash crittografica $H(\cdot)$ (SHA-256), e solo il digest risultante viene firmato tramite la chiave privata RSA. In nessuna fase dell'infrastruttura viene impiegata la firma RSA \textit{textbook} applicata direttamente al dato in chiaro.

\clearpage
\section{Il flusso del protocollo}
\label{sec:flusso_protocollo}
% Obiettivo della sezione: descrivere l'happy path cronologico, dalla preparazione
% dell'elezione fino alla registrazione della scheda.
%
% TODO: Inserire un diagramma di sequenza Elettore--SA--AE. Il diagramma deve mostrare:
% 1. TLS verso SA e autenticazione con credenziali istituzionali;
% 2. rilascio di C_p;
% 3. TLS verso AE;
% 4. invio di M = (C, C_p, eta, id_ele, sigma);
% 5. controlli AE;
% 6. rilascio di R;
% 7. cancellazione locale di PR_v.

Le Fasi dalla 0 alla 3 descrivono il percorso cronologico ideale dell'intero processo elettorale, dalla configurazione crittografica iniziale fino alla ricezione e registrazione sicura della scheda nell'urna pubblica. 

Il seguente diagramma di sequenza illustra ad alto livello le interazioni di rete tra i tre attori fondamentali (Elettore, Sistema di Autenticazione e Autorità Elettorale) durante le fasi attive di voto (Fasi 1-3). Ogni singola operazione riportata nel diagramma verrà analizzata tecnicamente nelle sottosezioni successive.

\vspace{1.5cm}
\begin{figure}[h!]
\centering
\begin{verbatim}
  Elettore                        SA                           AE
     |                             |                            |
     |-- 1. Handshake TLS -------->|                            |
     |      (Autenticazione con    |                            |
     |      credenziali d'Ateneo)  |                            |
     |                             |                            |
     |<-- 2. Rilascio di C_p ------|                            |
     |                             |                            |
     |-- 3. Handshake TLS ------------------------------------->|
     |                                                          |
     |   4. Invio del pacchetto di voto:                        |
     |      M = (C, C_p, \eta, id_{elezione}, \sigma)           |
     |--------------------------------------------------------->|
     |                                                          |
     |                                           5. [Esecuzione |
     |                                               Controlli] |
     |                                                          |
     |<--------------------------------------- 6. Rilascio di R |
     |                             |                            |
     | 7. [Cancellazione locale    |                            |
     |     della chiave PR_v]      |                            |
\end{verbatim}
\caption{Diagramma di sequenza del flusso di sottomissione del voto.}
\label{fig:diagramma_sequenza}
\end{figure}

\subsection*{Fase 0 --- Setup pre-elezione e distribuzione delle chiavi}
\label{subsec:fase_setup}
\label{subsec:wp2_fase_setup}

Prima dell'apertura delle urne, le entità di gestione provvedono alla configurazione crittografica del sistema e alla distribuzione sicura delle chiavi.

\subsection*{Generazione delle chiavi}
\label{subsubsec:wp2_generazione_chiavi}
% TODO: Descrivere la generazione di due coppie RSA distinte da parte dell'AE:
% (PU_ele, PR_ele) per la cifratura/decifratura delle schede e
% (PU_AE, PR_AE) per la firma di record, ricevute e documento finale.
%
% TODO: Motivare la separazione delle chiavi come scelta di least privilege.
%
% TODO: Descrivere la generazione della coppia (PU_SA, PR_SA) da parte del SA per
% la firma dei certificati pseudonimi.

L'Autorità Elettorale (AE) genera localmente due coppie di chiavi RSA distinte e tra loro indipendenti:
\begin{itemize}
    \item $(PU_{ele}, PR_{ele})$: coppia dedicata esclusivamente alla cifratura delle schede da parte degli elettori e alla successiva decifratura da parte dell'AE.
    \item $(PU_{AE}, PR_{AE})$: coppia dedicata esclusivamente all'apposizione delle firme digitali sui record dell'urna, sulle ricevute e sul documento di scrutinio finale.
\end{itemize}
L'isolamento di questi due domini operativi risponde rigorosamente al principio del minimo privilegio (\textit{least privilege}). In questo modo, l'eventuale compromissione della chiave di firma $PR_{AE}$ non implica in alcun modo la compromissione di $PR_{ele}$. Quest'ultima rimane in possesso esclusivo dell'AE, assicurando che la segretezza dei voti resti protetta per l'intera durata della votazione.

Parallelamente, il Sistema di Autenticazione (SA) genera la propria coppia di chiavi RSA $(PU_{SA}, PR_{SA})$, la cui chiave privata verrà impiegata unicamente per firmare digitalmente i certificati pseudonimi rilasciati agli elettori abilitati.

\subsection*{Distribuzione delle chiavi pubbliche tramite PKI}
\label{subsubsec:wp2_pki}
% TODO: Spiegare che PU_ele, PU_AE e PU_SA sono certificate dalla CA d'Ateneo tramite
% certificati X.509 e pubblicate prima dell'apertura delle urne.
%
% TODO: Collegare la stessa PKI anche all'autenticazione dei server nei canali TLS.

Al fine di garantire l'autenticità del materiale crittografico, il protocollo si appoggia all'Infrastruttura a Chiave Pubblica (PKI) istituzionale. Le tre chiavi pubbliche $PU_{ele}$, $PU_{AE}$ e $PU_{SA}$ vengono certificate dalla Certificate Authority (CA) dell'Ateneo tramite certificati digitali in formato X.509. Tali certificati vengono pubblicati nel \textit{repository} istituzionale e resi disponibili a tutti gli attori prima dell'inizio delle operazioni di voto.

Questa infrastruttura PKI svolge un duplice ruolo operativo: oltre a distribuire le chiavi crittografiche applicative, i medesimi certificati X.509 vengono impiegati per l'autenticazione dei server (SA e AE) durante la fase di instaurazione dei canali sicuri TLS. Ciò impedisce nativamente che un avversario di rete possa impersonare i server legittimi.

\subsection*{Blocco di testa dell'urna}
\label{subsubsec:wp2_blocco_testa}
% TODO: Presentare il record iniziale dell'urna:
% Params = { id_ele, PU_ele, N, lista_candidati, t_apertura, t_chiusura }.
%
% TODO: Specificare che Params viene firmato dall'AE e costituisce Record_0 della
% hash chain.

A conclusione della fase di setup, l'AE provvede a inizializzare l'urna pubblica. L'inizializzazione avviene mediante la pubblicazione di un blocco contenente i parametri pubblici e immutabili della consultazione elettorale:
$$Params = \{ id_{elezione}, PU_{ele}, N, \text{lista\_candidati}, t_{apertura}, t_{chiusura} \}$$
Questo blocco viene firmato digitalmente dall'AE tramite la propria chiave privata, producendo una firma applicata all'hash dei parametri ($Sign_{PR_{AE}}(H(Params))$). Tale struttura, pubblicata e verificabile da chiunque possieda $PU_{AE}$, costituisce il record di testa, denominato $Record_0$. Esso funge da ancoraggio crittografico di partenza per l'intera catena di hash (\textit{hash chain}) che comporrà l'urna pubblica.

\subsection*{Fase 1 --- Autenticazione ed emissione del certificato pseudonimo}
\label{subsec:fase_autenticazione}
\label{subsec:wp2_fase_autenticazione}

\subsection*{Generazione della coppia effimera lato elettore}
\label{subsubsec:wp2_chiave_effimera}
% TODO: Descrivere la generazione locale di (PU_v, PR_v), priva di informazioni
% identificative e valida solo per la sessione di voto.
%
% TODO: Precisare che PR_v non lascia il dispositivo dell'elettore e verrà cancellata
% dopo la ricezione della ricevuta R.

Prima di stabilire qualsiasi contatto di rete con le autorità, l'elettore genera localmente sul proprio dispositivo una coppia di chiavi RSA monouso $(PU_v, PR_v)$. Questa coppia costituisce l'identità crittografica pseudonima dell'elettore per la singola sessione di voto ed è intrinsecamente priva di qualsiasi informazione che possa ricondurre alla sua identità anagrafica. 

Per garantire l'assenza di prove crittografiche riutilizzabili a fini di coercizione, la chiave privata $PR_v$ non abbandona mai la memoria locale del dispositivo dell'elettore. La sua finestra di validità è deliberatamente ridotta al minimo: essa verrà infatti cancellata in modo sicuro e definitivo non appena la procedura di voto sarà completata, ovvero immediatamente dopo la ricezione e verifica della ricevuta $R$ da parte dell'Autorità Elettorale.

\subsection*{Autenticazione presso il SA}
\label{subsubsec:wp2_autenticazione_sa}
% TODO: Descrivere il canale TLS verso il SA e l'autenticazione con credenziali
% istituzionali.
%
% TODO: Indicare i due controlli del SA: appartenenza agli aventi diritto e assenza
% di un certificato pseudonimo già emesso per id_ele.

Terminata la fase di generazione locale, l'applicazione client dell'elettore stabilisce un canale sicuro TLS verso il Sistema di Autenticazione (SA), validandone il certificato X.509 emesso dalla CA d'Ateneo. All'interno di questo tunnel cifrato, l'elettore si autentica utilizzando le proprie credenziali istituzionali.

Il SA, avendo visibilità sull'identità reale dell'utente, esegue due controlli stringenti prima di procedere:
\begin{enumerate}
    \item Verifica l'appartenenza dell'utente all'elenco degli aventi diritto al voto per la specifica elezione in corso.
    \item Interroga il proprio registro interno (indicizzato per identità reale) per accertarsi che l'utente non abbia già richiesto e ottenuto un certificato pseudonimo per il medesimo $id_{elezione}$.
\end{enumerate}

\subsection*{Certificato pseudonimo}
\label{subsubsec:wp2_certificato_pseudonimo}
% TODO: Mostrare la struttura:
% C_p = { PU_v, id_ele, t_finestra, Sign_{PR_SA}(H(PU_v || id_ele || t_finestra)) }.
%
% TODO: Spiegare che C_p attesta l'abilitazione al voto senza includere l'identità
% reale dell'elettore.
%
% TODO: Motivare t_finestra come timestamp grossolano utile a ridurre la correlazione
% temporale tra emissione del certificato e arrivo della scheda all'AE.

Se i controlli danno esito positivo, il SA procede all'emissione del certificato pseudonimo $C_p$, strutturato come segue:
$$C_p = \{ PU_v, id_{elezione}, t_{finestra}, Sign_{PR_{SA}}(H(PU_v \parallel id_{elezione} \parallel t_{finestra})) \}$$

Questo certificato agisce come \textit{token} di abilitazione: attesta crittograficamente che la chiave pubblica effimera $PU_v$ appartiene a un elettore avente diritto per la votazione indicata, \textbf{senza includere in alcun campo l'identità reale dell'elettore}. Terminato il rilascio, il SA aggiorna il proprio registro interno (che non verrà mai condiviso con l'AE) e interrompe la comunicazione.

Il parametro $t_{finestra}$ inserito nel certificato costituisce un identificatore temporale a bassa risoluzione temporale, associato a un intero blocco orario. Assegnando il medesimo $t_{finestra}$ a numerosi certificati emessi a elettori diversi, si degrada volutamente la precisione del dato cronologico a disposizione dell'AE. Questa scelta progettuale riduce drasticamente la correlazione temporale tra il momento in cui l'elettore ottiene il certificato dal SA e l'istante in cui la scheda arriverà all'AE, proteggendo lo pseudo-anonimato contro tentativi di analisi del traffico (\textit{traffic analysis}).

\textbf{Proprietà realizzate:} [AUTH-1], [AUTH-2], [SEC-2] (condizionale all'Assunzione 1 di WP1: non-collusione SA/AE).

\subsection*{Fase 2 --- Costruzione e trasmissione della scheda}
\label{subsec:fase_costruzione_scheda}
\label{subsec:wp2_fase_costruzione_scheda}

\subsection*{Codifica scalare del voto}
\label{subsubsec:wp2_codifica_voto}
% TODO: Definire la preferenza come i in {1, ..., N}.
%
% TODO: Motivare la codifica scalare: RSA-OAEP non è omomorfico, quindi la codifica
% one-hot non dà vantaggi nel protocollo concreto. La validazione diventa il controllo
% 1 <= i <= N.

L'elettore seleziona la propria preferenza e la codifica in chiaro come un intero scalare $i \in \{1, \dots, N\}$. La scelta della codifica scalare è determinata direttamente dallo schema crittografico adottato: poiché RSA-OAEP non possiede proprietà omomorfiche, l'impiego di una codifica \textit{one-hot} non offrirebbe alcun beneficio per le operazioni di conteggio. L'Autorità Elettorale dovrà necessariamente decifrare le singole schede in modo individuale per calcolare il risultato aggregato. Di conseguenza, la validazione strutturale del voto si riduce a un semplice controllo di intervallo $1 \le i \le N$ eseguito in fase di scrutinio.

\subsection*{Cifratura con RSA-OAEP}
\label{subsubsec:wp2_rsa_oaep}
% TODO: Presentare C = RSA-OAEP-Enc(PU_ele, i; r).
%
% TODO: Spiegare perché OAEP è necessario su un dominio di messaggi piccolo e noto:
% senza randomizzazione, un avversario potrebbe cifrare tutte le opzioni e confrontarle.
%
% TODO: Chiarire che r resta noto solo all'elettore e contribuisce alla non rivelazione
% del voto rispetto a terzi che osservano il ciphertext.

Il voto in chiaro viene cifrato utilizzando la chiave pubblica dell'elezione, ottenendo il ciphertext:
$$C = \text{RSA-OAEP-Enc}(PU_{ele}, i; r)$$
L'utilizzo del padding OAEP e del fattore di randomizzazione $r$ è un requisito indispensabile per garantire la sicurezza semantica (CPA-security). Poiché il dominio dei messaggi possibili è molto piccolo e noto a priori (le $N$ opzioni di voto), l'applicazione di un algoritmo deterministico permetterebbe a un avversario di pre-calcolare i crittogrammi di tutte le liste e confrontarli con la scheda, identificando istantaneamente il voto. 

Il fattore $r$ rende ogni operazione di cifratura crittograficamente unica e indistinguibile. Esso viene generato localmente e rimane noto esclusivamente all'elettore. La sua segretezza concorre in modo sostanziale alla protezione della preferenza: un soggetto terzo che osservi il ciphertext $C$ nell'urna non può in alcun modo risalire a $i$ senza conoscere $r$.

\subsection*{Nonce eta e fingerprint della scheda}
\label{subsubsec:wp2_nonce_fingerprint}
% TODO: Descrivere la generazione di eta a 256 bit.
%
% TODO: Aggiornare la formula della verifica individuale:
% K = H(C), F = H(K || eta).
%
% TODO: Spiegare il doppio ruolo di eta:
% (i) segreto necessario a individuare il proprio F_j senza produrre una ricevuta
%     rivelante;
% (ii) identificatore di freschezza complementare al controllo di unicità di PU_v.
%
% TODO: Usare la terminologia "assenza di ricevuta rivelante", non receipt-freeness
% piena.

L'elettore genera localmente un valore casuale monouso a 256 bit:
$$\eta \leftarrow\$ \{0,1\}^{256}$$
Questo parametro (\textit{nonce}) riveste un duplice ruolo strutturale all'interno del protocollo:
\begin{enumerate}
    \item \textbf{Segreto per la verifica individuale:} partendo dal \textit{commitment} pubblico al ciphertext calcolato come $K = H(C)$, il \textit{fingerprint} univoco della scheda viene derivato come $F = H(K \parallel \eta)$. Essendo $\eta$ noto esclusivamente all'elettore, nessun terzo può calcolare questo hash e individuare a quale utente corrisponda una determinata scheda nell'urna. Ciò abilita il meccanismo di verifica della presenza del voto garantendo l'\textbf{assenza di ricevuta rivelante}, poiché l'elettore non dispone di alcun elemento esibibile per dimostrare il proprio voto a un potenziale coercitore.
    \item \textbf{Parametro anti-replay:} agisce come misura complementare. Sebbene l'unicità del voto sia garantita in prima istanza dal controllo sulla chiave effimera $PU_v$, il nonce $\eta$ fornisce un ulteriore livello di protezione crittografica, verificabile in modo indipendente e ancor prima di validare lo stato del certificato pseudonimo.
\end{enumerate}

\subsection*{Pacchetto di voto ed Encrypt-then-Sign}
\label{subsubsec:wp2_encrypt_then_sign}
% TODO: Presentare:
% sigma = Sign_{PR_v}(H(C || C_p || eta || id_ele));
% M = (C, C_p, eta, id_ele, sigma).
%
% TODO: Spiegare che la firma copre il ciphertext C, non il plaintext i, e rende
% rilevabile qualsiasi sostituzione del ciphertext in transito.
%
% TODO: Chiarire che con Sign si intende RSA hash-and-sign, non RSA textbook.
%
% TODO: Ribadire che eta viaggia nel pacchetto M, ma è protetto in transito da TLS.

Il processo di impacchettamento si conclude seguendo rigorosamente il paradigma \textit{Encrypt-then-Sign}. L'elettore firma l'intero set di dati utilizzando la propria chiave privata effimera e trasmette il pacchetto $M$:
$$\sigma = \text{Sign}_{PR_v}(H(C \parallel C_p \parallel \eta \parallel id_{elezione}))$$
$$M = (C, C_p, \eta, id_{elezione}, \sigma)$$

In aderenza alle specifiche di sicurezza, l'operazione $\text{Sign}$ impiega unicamente lo schema RSA \textit{hash-and-sign}, impedendo l'uso di firme RSA \textit{textbook}.

È cruciale sottolineare che la firma $\sigma$ copre l'hash del ciphertext $C$, e non il voto in chiaro $i$. Questo meccanismo rileva qualsiasi tentativo di manipolazione in transito: un attaccante \textit{Man-in-the-Middle} che provasse a sostituire $C$ con una propria scheda cifrata invaliderebbe matematicamente $\sigma$. 

Infine, si ribadisce che $\eta$, $C_p$ e la firma viaggiano in chiaro all'interno della struttura logica del pacchetto $M$, affinché l'AE possa estrarli ed elaborare i controlli previsti. Tuttavia, la loro confidenzialità lungo la trasmissione in rete è integralmente protetta dal tunnel TLS sottostante.

\textbf{Proprietà realizzate:} [INT-1], [UNIQ-2].

\subsection*{Fase 3 --- Ricezione, validazione e rilascio della ricevuta}
\label{subsec:fase_ricezione}
\label{subsec:wp2_fase_ricezione}

\subsection*{Controlli dell'AE}
\label{subsubsec:wp2_controlli_ae}
% TODO: Inserire pseudocodice con i quattro controlli sequenziali:
% 1. verifica firma SA su C_p;
% 2. unicità di PU_v nel Registro_Effimeri_AE;
% 3. anti-replay su eta nel Registro_Nonce_AE;
% 4. verifica sigma tramite PU_v.
%
% TODO: Specificare che il fallimento di uno qualunque dei controlli causa il rigetto
% immediato del pacchetto.

L'Autorità Elettorale riceve via TLS il pacchetto $M = (C, C_p, \eta, id_{elezione}, \sigma)$ ed esegue una validazione crittografica e procedurale prima di accettare il voto. I controlli avvengono in modo strettamente sequenziale:
\begin{enumerate}
    \item \textbf{Verifica Autorizzazione:} validazione della firma del SA posta sul certificato $C_p$, per accertare che l'utente sia un avente diritto abilitato.
    \item \textbf{Verifica Unicità:} controllo che la chiave pubblica $PU_v$ contenuta in $C_p$ non sia già presente nel \texttt{Registro\_Effimeri\_AE}. Questo impedisce a un utente di votare due volte con lo stesso certificato.
    \item \textbf{Verifica Anti-replay:} controllo che il nonce $\eta$ non sia mai stato registrato nel \texttt{Registro\_Nonce\_AE}.
    \item \textbf{Verifica Integrità:} validazione della firma effimera $\sigma$ applicata sull'hash del pacchetto, impiegando la chiave pubblica $PU_v$ appena validata.
\end{enumerate}
L'architettura impone una logica binaria di accettazione: il fallimento di anche uno solo di questi controlli causa il rigetto immediato della richiesta. Il pacchetto anomalo viene scartato senza essere inserito nell'urna e la connessione interrotta con un messaggio di errore.

\subsection*{Commitment, fingerprint e ricevuta}
\label{subsubsec:wp2_commitment_ricevuta}
% TODO: Aggiornare il passo di accettazione:
% K_j = H(C_j);
% F_j = H(K_j || eta_j);
% R = Sign_{PR_AE}(H(K_j || F_j || j || id_ele)).
%
% TODO: Specificare che j è l'indice progressivo assegnato dall'AE e che R viene
% inviata all'elettore via TLS.

Se il pacchetto supera positivamente tutti i controlli, l'AE procede alla sua accettazione. Al pacchetto viene assegnato un indice identificativo progressivo $j$. L'autorità calcola quindi due valori fondamentali:
$$K_j = H(C_j) \quad \text{(Commitment pubblico)}$$
$$F_j = H(K_j \parallel \eta_j) \quad \text{(Fingerprint)}$$

Per concludere la transazione, l'AE genera la ricevuta crittografica $R$ firmando digitalmente i parametri essenziali dell'operazione:
$$R = \text{Sign}_{PR_{AE}}(H(K_j \parallel F_j \parallel j \parallel id_{elezione}))$$
La ricevuta $R$ viene immediatamente inviata all'elettore attraverso il canale TLS ancora aperto, a testimonianza dell'avvenuta ricezione e validazione.

\subsection*{Separazione tra archivio interno e urna pubblica}
\label{subsubsec:wp2_archivio_vs_urna}
% TODO: Inserire una tabella o blocco descrittivo:
% Archivio interno AE: C_j, C_p, eta_j, sigma, PU_v, stato della scheda.
% Urna pubblica: j, K_j, F_j, H(Record_{j-1}), Sign_AE(...).
%
% TODO: Chiarire che C_j serve allo scrutinio ma non viene pubblicato; K_j è il
% commitment pubblico al ciphertext.

Per preservare la confidenzialità del voto senza rinunciare alla verificabilità, l'AE implementa una netta separazione tra i dati conservati privatamente e quelli pubblicati. 
Il ciphertext $C_j$ è indispensabile per la successiva fase di scrutinio, ma esporlo pubblicamente renderebbe il voto vulnerabile ad attacchi crittografici futuri. Pertanto, l'AE pubblica esclusivamente il commitment $K_j$, il quale vincola crittograficamente l'AE al ciphertext ricevuto, senza però rivelarlo in chiaro.

La Tabella \ref{tab:separazione_dati} illustra la ripartizione informativa:

\begin{table}[H]
\centering
\renewcommand{\arraystretch}{1.3}
\begin{tabular}{|p{7cm}|p{7.5cm}|}
\hline
\textbf{Archivio Interno AE (Privato)} & \textbf{Urna Pubblica (Record pubblicato)} \\ \hline
Scheda cifrata originale ($C_j$). & Indice progressivo $j$. \\ \hline
Certificato pseudonimo $C_p$ (e chiave $PU_v$). & \textit{Commitment} al ciphertext ($K_j = H(C_j)$). \\ \hline
Nonce $\eta_j$ e firma effimera $\sigma$. & \textit{Fingerprint} della scheda ($F_j$). \\ \hline
Stato di elaborazione della scheda. & Hash pointer al record precedente e firma dell'AE. \\ \hline
\end{tabular}
\caption{Separazione dei dati tra archivio privato e registro pubblico.}
\label{tab:separazione_dati}
\end{table}

\subsection*{Cancellazione della chiave effimera}
\label{subsubsec:wp2_cancellazione_prv}
% TODO: Descrivere la cancellazione immediata di PR_v dopo la verifica della ricevuta.
%
% TODO: Collegare la cancellazione alla mitigazione della coercizione retroattiva:
% l'elettore non può produrre nuove firme che dimostrino la paternità del pacchetto.
%
% TODO: Specificare che l'elettore conserva quanto serve alla verifica individuale:
% C, eta, R e l'indice j ricavabile da R; K_j è ricalcolabile come H(C) o recuperabile
% dal record pubblico.

Non appena il client dell'elettore riceve e verifica la correttezza matematica della ricevuta $R$, l'applicazione procede alla cancellazione istantanea e irrecuperabile della chiave privata $PR_v$ dal dispositivo locale.

Questo passaggio non è una semplice routine di pulizia della memoria, ma una misura primaria per la mitigazione della \textbf{coercizione retroattiva}. Distruggendo la chiave utilizzata per firmare il pacchetto, l'elettore si priva intenzionalmente della capacità crittografica di provare la paternità della firma $\sigma$. Qualora un coercitore richiedesse in un secondo momento di dimostrare l'associazione tra la propria identità e un record nell'urna, l'elettore non potrebbe fornire alcuna prova incontestabile. 

Ai fini della verifica individuale VER-1, l'elettore conserva localmente soltanto il ciphertext $C$, il nonce $\eta$, la ricevuta $R$ e l'indice progressivo $j$ (deducibile da $R$). Il commitment $K_j$ non necessita di essere salvato, in quanto l'elettore può ricalcolarlo autonomamente applicando l'hash a $C$, oppure può semplicemente recuperarlo dal record pubblico dell'urna.

\textbf{Proprietà realizzate:} [INT-2], [UNIQ-1], [SEC-3].

\clearpage
\section{Architettura dell'urna pubblica: struttura ibrida}
\label{sec:architettura_urna}
% Obiettivo della sezione: descrivere come vengono pubblicati e verificati i record
% dell'urna.
%
% TODO: Introdurre la combinazione hash chain + Merkle Tree, richiamando hash pointer,
% tamper-evident log e proof of membership.
%
% TODO: Separare chiaramente i ruoli: hash chain come log storico; Merkle Tree come
% indice efficiente per la verifica individuale.

\subsection*{Append-only log basato su hash chain}
\label{subsec:hash_chain}
\label{subsec:wp2_hash_chain}
% TODO: Aggiornare la struttura del record pubblico:
% Record_j = { j, K_j, F_j, H(Record_{j-1}), Sign_{PR_AE}(H(j || K_j || F_j || H(Record_{j-1}))) }.
%
% TODO: Spiegare che Record_0 è il blocco di testa Params.
%
% TODO: Descrivere perché hash pointer e firme rendono rilevabili modifiche retroattive.

\subsection*{Ruolo del commitment pubblico \(K_j\)}
\label{subsubsec:wp2_ruolo_kj}
% TODO: Spiegare che K_j = H(C_j) non rivela C_j, ma impegna pubblicamente l'AE al
% ciphertext ricevuto.
%
% TODO: Specificare che K_j abilita controlli procedurali o audit: un revisore con
% accesso a C_j può verificare K_j = H(C_j).
%
% TODO: Non presentare K_j come soluzione a VER-2 piena; è un rafforzamento strutturale.

\subsection*{Pubblicazione per finestre}
\label{subsubsec:wp2_pubblicazione_finestre}
% TODO: Dichiarare che i record non vengono pubblicati in tempo reale uno per uno, ma
% per finestre temporali coerenti con t_finestra.
%
% TODO: Motivare la scelta come mitigazione della correlazione temporale SA--AE.

\subsection*{Indice efficiente basato su Merkle Tree}
\label{subsec:merkle_tree}
\label{subsec:wp2_merkle_tree}
% TODO: Spiegare che i fingerprint F_j sono indicizzati in un Merkle Tree.
%
% TODO: Specificare che la Merkle Root viene pubblicata alla fine di ogni finestra e
% definitivamente alla chiusura delle urne.
%
% TODO: Distinguere Verificatore Universale (scarica la hash chain) ed elettore
% individuale (scarica solo la proof di membership relativa a F_j).

\subsection*{Scelta progettuale: architettura centralizzata vs blockchain}
\label{subsec:wp2_centralizzata_vs_blockchain}
% TODO: Spiegare che il design usa primitive tipiche della blockchain, ma senza adottare
% una blockchain completa.
%
% TODO: Argomentare i quattro motivi:
% 1. presenza di una trust authority istituzionale;
% 2. efficienza e latenza;
% 3. la blockchain non risolve lo scrutinio black-box;
% 4. lo smart contract non aggiunge valore crittografico rispetto alla hash chain firmata.
%
% TODO: Lasciare la blockchain come possibile direzione futura in scenari trustless o
% con decifratura distribuita, ma non come scelta del WP2.

\section{Chiusura delle urne, scrutinio e pubblicazione}
\label{sec:scrutinio_pubblicazione}

\subsection*{Chiusura e congelamento dell'urna}
\label{subsec:wp2_chiusura}
% TODO: Descrivere il record di chiusura firmato dall'AE, contenente Merkle Root finale
% e numero totale di schede.
%
% TODO: Specificare che dopo t_chiusura non vengono accettati nuovi record.

\subsection*{Decifratura e validazione delle schede}
\label{subsec:wp2_decifratura_validazione}
% TODO: Specificare che l'AE recupera C_j dall'archivio interno, non dall'urna pubblica.
%
% TODO: Presentare i_j = RSA-OAEP-Dec(PR_ele, C_j).
%
% TODO: Applicare il controllo di conformità 1 <= i_j <= N.
%
% TODO: Le schede non conformi vengono marcate come invalide con record firmato e non
% eliminate dall'urna.

\subsection*{Calcolo e pubblicazione del risultato}
\label{subsec:wp2_calcolo_risultato}
% TODO: Presentare Risultato[k] = |{j : i_j = k}|.
%
% TODO: Descrivere Doc_finale con Risultato, k_totale, k_conformi, k_invalide,
% Merkle_Root_finale e firma AE.
%
% TODO: Specificare che i singoli plaintext i_j non vengono pubblicati.

\subsection*{Dichiarazione del vincolo strutturale \texttt{[SEC-1]}}
\label{subsec:limite_sec1}
\label{subsec:wp2_limite_sec1}
% TODO: Dichiarare che SEC-1 piena richiederebbe che nemmeno l'AE possa leggere le
% singole schede, ma RSA-OAEP non è omomorfico e quindi l'AE deve decifrare C_j.
%
% TODO: Citare come strumenti non trattati: cifratura omomorfica additiva, decifratura
% a soglia, mix-net verificabile.
%
% TODO: Specificare quanto è realizzato: segretezza in transito tramite TLS, segretezza
% in conservazione pubblica perché l'urna contiene K_j e F_j ma non C_j, accesso ai
% voti in chiaro limitato alla fase di scrutinio e all'AE.

\section{Meccanismi di verifica}
\label{sec:meccanismi_verifica}

\subsection*{Verifica individuale \texttt{[VER-1]}}
\label{subsec:verifica_individuale}
\label{subsec:wp2_verifica_individuale}
% TODO: Descrivere i passi dell'elettore:
% 1. calcola K = H(C);
% 2. calcola F = H(K || eta);
% 3. scarica la Merkle Proof per F e verifica la root finale;
% 4. scarica Record_j e verifica K_j = K.
%
% TODO: Spiegare che la ricevuta R lega K_j, F_j, j e id_ele tramite firma AE.
%
% TODO: Chiarire che la proprietà è "assenza di ricevuta rivelante rilasciata dal
% sistema", non impossibilità assoluta di costruire prove esterne.
%
% TODO: Specificare il caso di scheda scartata: l'elettore può esibire R per attivare
% una verifica procedurale.

\subsection*{Verifica universale \texttt{[VER-2]}}
\label{subsec:verifica_universale}
\label{subsec:wp2_verifica_universale}
% TODO: Elencare i controlli pubblici:
% - consistenza della Merkle Root con i fingerprint F_j;
% - firme valide su tutti i record;
% - hash pointer corretti;
% - coerenza k_conformi + k_invalide = k_totale;
% - firma valida su Doc_finale.
%
% TODO: Aggiungere il caso auditor: se un revisore ottiene accesso ai ciphertext C_j,
% può verificare K_j = H(C_j) e quindi la coerenza tra archivio interno e urna pubblica.

\subsection*{Dichiarazione del vincolo strutturale \texttt{[VER-2]}}
\label{subsec:limite_ver2}
\label{subsec:wp2_limite_ver2}
% TODO: Dichiarare che il protocollo realizza una verificabilità universale strutturale,
% non una verifica matematica piena del conteggio.
%
% TODO: Spiegare che resta non verificabile pubblicamente se l'AE abbia decifrato e
% conteggiato correttamente ogni C_j.
%
% TODO: Citare strumenti non trattati che servirebbero per VER-2 piena: ZKP di corretta
% decifratura, tally omomorfico, mix-net verificabile.
%
% TODO: Evidenziare le due mitigazioni disponibili: verifica individuale con R e audit
% procedurale basato su K_j = H(C_j), se i ciphertext vengono resi disponibili ai revisori.

\section{Riepilogo funzionale delle proprietà di sicurezza}
\label{sec:riepilogo_proprieta}
% TODO: Inserire una tabella di tracciabilità con colonne:
% Proprietà | Stato realizzato | Meccanismo del protocollo | Assunzione/Vincolo.
%
% TODO: Aggiornare le righe alla versione con K_j:
% AUTH-1, AUTH-2, UNIQ-1, UNIQ-2, INT-1, INT-2, INT-3,
% SEC-1, SEC-2, SEC-3, VER-1, VER-2.
%
% TODO: Usare per SEC-3 la formula "assenza di ricevuta rivelante" e per VER-2 la
% formula "strutturale".
%
% TODO: Specificare che la tabella non è ancora analisi di sicurezza: la resilienza
% effettiva rispetto agli attaccanti del WP1 sarà discussa in WP3.

\section{Analisi critica delle scelte progettuali}
\label{sec:scelte_progettuali}

\subsection*{Modello one-shot irrevocabile}
\label{subsec:compromesso_irrevocabilita}
\label{subsec:wp2_one_shot_irrevocabile}
% TODO: Dichiarare che la scelta di non revoca/non sostituzione del voto è definitiva.
%
% TODO: Argomentare i tre conflitti della sostituzione:
% 1. richiederebbe un identificativo persistente;
% 2. entrerebbe in conflitto con la cancellazione di PR_v;
% 3. complicherebbe VER-1 e VER-2.
%
% TODO: Dichiarare il rischio residuo di coercizione pre-invio come accettato nel modello.

\subsection*{Architettura centralizzata vs blockchain}
\label{subsec:compromesso_blockchain}
\label{subsec:wp2_esclusione_blockchain}
% TODO: Rinviare alla Sezione~\ref{subsec:wp2_centralizzata_vs_blockchain} e riassumere
% le ragioni della scelta: trust authority, efficienza, scrutinio black-box, scarso valore
% aggiunto dello smart contract.

\section{Strumenti del corso utilizzati e considerati}
\label{sec:strumenti_corso}

\subsection*{Strumenti adottati}
\label{subsec:strumenti_adottati}
\label{subsec:wp2_strumenti_adottati}
% TODO: Elencare e mappare alle componenti del protocollo:
% - RSA;
% - RSA-OAEP;
% - firma digitale RSA hash-and-sign;
% - SHA-256;
% - Merkle Tree;
% - hash chain / tamper-evident log;
% - PKI e certificati X.509;
% - TLS;
% - Encrypt-then-Sign.
%
% TODO: Nella voce firma digitale, ribadire che Sign non indica firma RSA textbook.
%
% TODO: Nella voce TLS, ribadire che TLS protegge i contenuti ma non elimina del tutto
% i canali laterali di rete; la mitigazione è completata da t_finestra e pubblicazione
% per finestre.

\subsection*{Strumenti considerati e non adottati}
\label{subsec:strumenti_non_adottati}
\label{subsec:wp2_strumenti_non_adottati}
% TODO: Motivare l'esclusione di:
% - AES e cifratura simmetrica applicativa dedicata;
% - KDC e Needham-Schroeder;
% - Diffie-Hellman applicativo diretto;
% - blockchain e smart contract.
%
% TODO: Non includere come alternative operative blind signatures, omomorfismo, ZKP,
% mix-net o threshold decryption, perché non sono strumenti trattati nel corso. Possono
% essere citati solo nelle sezioni sui vincoli strutturali come strumenti esterni che
% servirebbero per ottenere SEC-1 o VER-2 in forma piena.
